% ===== PRÉAMBULE CANONIQUE — à coller VERBATIM en tête de CHAQUE .tex =====
\documentclass[11pt,a4paper]{article}
\usepackage[utf8]{inputenc}\usepackage[T1]{fontenc}\usepackage{lmodern}
\usepackage[french]{babel}
\usepackage{amsmath,amssymb,mathtools}\usepackage{icomma}
\usepackage[version=4]{mhchem}          % \ce{...} : équations & formules
\usepackage{siunitx}\sisetup{locale=FR} % \qty{}{}, \si{}, \ang{}
\usepackage{chemfig}                    % \chemfig{...} : structures développées
\usepackage{xcolor}\usepackage{colortbl}
\usepackage{tikz}\usepackage{pgfplots}\pgfplotsset{compat=1.18}
\usepackage[most]{tcolorbox}\usepackage{enumitem}\usepackage{array,booktabs}
\usepackage[a4paper,margin=2.2cm]{geometry}
\usetikzlibrary{arrows.meta,calc,angles,quotes,positioning,patterns,backgrounds,shapes.geometric}
\definecolor{primary}{RGB}{222,49,99}\definecolor{primarydark}{RGB}{150,22,55}
\definecolor{defcol}{RGB}{0,114,178}\definecolor{methcol}{RGB}{52,92,148}
\definecolor{excol}{RGB}{0,135,90}\definecolor{attncol}{RGB}{200,40,55}
\tcbset{boxbase/.style={enhanced,breakable,boxrule=0.9pt,arc=2.5pt,
  left=3mm,right=3mm,top=2.3mm,bottom=2.3mm,fonttitle=\bfseries,coltitle=white}}
% --- boîtes TUTORIEL (noms EXACTS, ne pas renommer) ---
\newtcolorbox{cle}[1][]{boxbase,colback=defcol!6,colframe=defcol,
  title={\textsc{À retenir}\ifx\relax#1\relax\relax\else\textmd{ : \itshape #1}\fi}}
\newtcolorbox{methode}[1][]{boxbase,colback=methcol!7,colframe=methcol,sharp corners,
  title={\textsc{Méthode}\ifx\relax#1\relax\relax\else\textmd{ --- \itshape #1}\fi}}
\newtcolorbox{exemplebox}[1][]{boxbase,colback=excol!5,colframe=excol,
  title={\textsc{Exemple}\ifx\relax#1\relax\relax\else\textmd{ : \itshape #1}\fi}}
\newtcolorbox{piege}[1][]{boxbase,colback=attncol!6,colframe=attncol,
  title={\textsc{Piège}\ifx\relax#1\relax\relax\else\textmd{ : \itshape #1}\fi}}
\newtcolorbox{remarque}[1][]{boxbase,colback=black!4,colframe=black!45,
  title={\textsc{Remarque}\ifx\relax#1\relax\relax\else\textmd{ : \itshape #1}\fi}}
% --- boîtes EXERCICE / CORRIGÉ (noms EXACTS, auto-comptées) ---
\newtcolorbox[auto counter]{exo}[1][]{boxbase,colback=primary!4,colframe=primary,
  title={\textsc{Exercice}~\thetcbcounter\ifx\relax#1\relax\relax\else\textmd{ --- \itshape #1}\fi}}
\newtcolorbox[auto counter]{corrige}[1][]{boxbase,colback=excol!4,colframe=excol,
  title={\textsc{Corrigé}~\thetcbcounter\ifx\relax#1\relax\relax\else\textmd{ --- \itshape #1}\fi}}
\newcommand{\R}{\mathbb{R}}\newcommand{\N}{\mathbb{N}}\newcommand{\Z}{\mathbb{Z}}\newcommand{\ds}{\displaystyle}
% (un \usepackage{fancyhdr}+en-tête est possible ; il est ignoré côté web)
% =========================================================================
\begin{document}
\begin{exo}[y a-t-il résonance ?]
La résonance apparaît dès qu'une double liaison est conjuguée à un doublet libre ou à une autre
double liaison. Pour chaque espèce, indique si elle présente une résonance (oui / non).
\begin{enumerate}[label=\alph*)]
  \item \ce{CO3^2-}
  \item \ce{CH4}
  \item \ce{O3}
  \item \ce{H2O}
  \item \ce{NO3-}
\end{enumerate}
\end{exo}

\begin{exo}[combien de formes limites équivalentes ?]
Donne le nombre de formes limites \emph{équivalentes} de chaque espèce.
\begin{enumerate}[label=\alph*)]
  \item \ce{O3}
  \item \ce{CO3^2-}
  \item \ce{NO3-}
  \item \ce{NO2-}
  \item \ce{SO2}
\end{enumerate}
\end{exo}

\begin{exo}[qu'est-ce qui bouge ?]
Vrai ou faux ? Pour chaque affirmation sur la résonance.
\begin{enumerate}[label=\alph*)]
  \item Les noyaux (atomes) changent de position d'une forme à l'autre.
  \item Les doublets $\pi$ et les doublets libres se délocalisent.
  \item La géométrie de la molécule change d'une forme à l'autre.
  \item On relie les formes limites par la flèche $\leftrightarrow$.
\end{enumerate}
\end{exo}

\begin{exo}[$\leftrightarrow$ ou $\rightleftharpoons$ ?]
Quel symbole convient à chaque situation ?
\begin{enumerate}[label=\alph*)]
  \item Deux formes limites de l'ozone.
  \item Un équilibre chimique entre réactifs et produits.
  \item Les trois structures de résonance du carbonate.
\end{enumerate}
\end{exo}

\begin{exo}[conséquence sur les liaisons]
Dans l'ion carbonate \ce{CO3^2-}, la résonance délocalise la double liaison sur les trois oxygènes.
\begin{enumerate}[label=\alph*)]
  \item Les trois liaisons \ce{C-O} réelles sont-elles identiques ou différentes ?
  \item Leur longueur est-elle celle d'une simple, d'une double, ou \emph{intermédiaire} ?
  \item Sur quels atomes la charge $-2$ est-elle répartie ?
\end{enumerate}
\end{exo}

\end{document}
